\chapter{Experimental Evidence}
\label{chap:experimental}

\section{Skein Decomposition and Hadronic Mass}
\label{sec:skein-hadronic-mass}

\begin{remark}[Skein Decomposition and Topological Mass Ratios]
  \label{rem:skein-decomposition}
This remark formalizes the refinement of the particle-topology mapping by shifting from torus knots to a framework grounded in \textbf{hyperbolic volumes} and \textbf{Skein-deformed residues}. By treating mass as the geometric resistance of the 1-jet bundle prolongation, the $q$-deformed Skein relation provides a rigorous "Mass Sum Rule" for composite hadronic states.

\textbf{Topological Ansatz and Field Character.}
The following table summarizes the axiomatic assignments where particle identities are mapped to specific knot types ($K$) and their associated masses $\mathcal{M}(K)$ are governed by the local substrate field $q(p)$.

$$
\begin{array}{|l|l|l|}
\hline
\textbf{Particle / State} & \textbf{Topological Identity} & \textbf{QOU Geometric Role} \\ \hline
\text{Photon} & \text{The Point (Empty Knot)} & \text{Background flux on the exceptional divisor} \\ \hline
\text{Neutrino} & \text{Unknot } (\mathcal{O}) & \text{Baseline vacuum drag; trivial braiding} \\ \hline
\text{Electron} & \text{Trefoil } (3_1) & \text{Simplest stable vortex; } deg~\alpha_V = 1 \text{ generation} \\ \hline
\text{Muon} & \text{Figure-Eight } (4_1) & \text{Hyperbolic volume residue; } deg~\alpha_V = 2 \text{ generation} \\ \hline
\text{Nucleon} & \text{Borromean Rings } (6_2^3) & \text{Color-exchange entanglement of three tubes} \\ \hline
\text{Strong Force} & \text{Hopf Link } (\mathcal{L}_{Hopf}) & \text{Internal binding energy; invisible to } \mathbb{P}^1(\mathbb{C}) \text{ base} \\ \hline
\end{array}
$$
\end{remark}

\textbf{Hyperbolic volumes and the topological mass ansatz.}
The key idea behind the following two propositions is that inertial mass
originates in the topology of knot complements. Every particle in the QOU
framework is identified with a knot or link $K$ (see the table in
\hyperref[rem:skein-decomposition]{Remark}). To each such knot one associates the
\textbf{hyperbolic volume} $\mathrm{Vol}(K)$: the volume of the complete
hyperbolic metric on the complement $S^3 \setminus K$, computed via the
Bloch--Wigner dilogarithm. For torus knots (e.g. the trefoil $3_1$) this
volume vanishes; for hyperbolic knots it is a positive topological invariant.
The figure-eight knot $4_1$, which the framework identifies with the muon,
has
\[
\mathrm{Vol}(4_1) = 3\,\mathrm{Cl}_2(\pi/3)
    \approx 2.02988\,321\,281\,931,
\]
where $\mathrm{Cl}_2$ is the Clausen function (OEIS A091518; SnapPy-verified).

The \textbf{topological mass conjecture} (Conjecture 4.25) asserts that the
mass $m(K)$ of a particle with knot type $K$ is
\[
m(K) \;=\; \frac{\mathrm{Vol}(K)}{\hbar_q^{\,2}},
  \qquad
  \hbar_q \;:=\; 1 - q^{-1},
\]
where $q > 1$ is the substrate parameter
(\hyperref[def:substrate-parameter]{Definition}) and $\hbar_q$ is the quantum
Planck constant (\hyperref[def:quantum-planck-constant]{Definition}). For
\emph{composite} states the mass functional $\mathcal{M}$ is \emph{not}
additive; instead it satisfies a $q$-deformed Skein recursion at every
crossing in the color-exchange diagram, generalising the Jones-polynomial
Skein relation to a mass-valued setting.

\begin{proposition}[$q$-Deformed Skein Relation]
  \label{prop:skein-mass-relation}
  \lean{QOU.HadronicMass.SkeinMassRelation}
  \uses{rem:skein-decomposition, def:substrate-parameter, def:quantum-planck-constant}
The mass functional $\mathcal{M}(H)$ for a composite hadron $H$ is not additive but obeys the quantum Skein relation at every crossing in the color-exchange space:
\[
q \,\mathcal{M}(L_+) - q^{-1}\,\mathcal{M}(L_-) = (q^{1/2} - q^{-1/2})\,\mathcal{M}(L_0).$$
Evaluation of the lepton mass ratio $m_{\mu}/m_{e}$ against **CODATA** recommended values yields a precise derivation for the local substrate parameter $q$:
\]
\frac{m_{\mu}}{m_{e}} = \frac{\mathrm{Vol}(4_1)}{\hbar_q^2}
  \approx \frac{2.02988}{(1 - q^{-1})^2},
\qquad 206.768 \implies q \approx 1.1097.
\[
---

The electron is identified with the trefoil $3_1$, a torus knot with
$\mathrm{Vol}(3_1) = 0$. Its mass therefore sits at the base of the
topological tower: it is the unit against which all other masses are measured.
The muon is identified with the figure-eight knot $4_1$, whose complement
carries a complete hyperbolic structure of volume
$\mathrm{Vol}(4_1) \approx 2.02988$.

Applying the topological mass conjecture to the muon--electron mass
*ratio* eliminates the overall mass scale and gives
\]
  \frac{m_\mu}{m_e}
  \;=\;
  \frac{\mathrm{Vol}(4_1) / \hbar_q^{\,2}}{m_e / m_e}
  \;=\;
  \frac{\mathrm{Vol}(4_1)}{\hbar_q^{\,2}}.
\[
Substituting $\hbar_q = 1 - q^{-1}$ and the CODATA value
$m_\mu / m_e = 206.768\,2827(46)$, one solves for $\hbar_q$:
\]
  \hbar_q^{\,2}
  = \frac{\mathrm{Vol}(4_1)}{m_\mu / m_e}
  = \frac{2.02988}{206.768}
  \approx 0.009\,817,
  \qquad
  \hbar_q \approx 0.09908.
\[
Inverting $\hbar_q = 1 - 1/q$ yields
$1/q = 1 - 0.09908 = 0.90092$, so
\]
  q \;\approx\; 1.1097.
\[
The Skein relation itself is the mass-functional analogue of the Kauffman
bracket recursion. At each crossing in the color-exchange diagram of a
hadron, the three resolutions $L_+$, $L_-$, $L_0$ satisfy
$q\,\mathcal{M}(L_+) - q^{-1}\mathcal{M}(L_-) =
(q^{1/2}-q^{-1/2})\,\mathcal{M}(L_0)$. This relation is inherited from the
$q$-deformed $\mathrm{SU}(2)$ Chern--Simons functional on the
$1$-jet bundle prolongation of the CR manifold
(see [Chapter](#chap:path-integrals-braiding)).
\]
\end{proposition}
\begin{proof}
The electron is identified with the trefoil $3_1$, a torus knot with
$\mathrm{Vol}(3_1) = 0$. Its mass therefore sits at the base of the
topological tower: it is the unit against which all other masses are measured.
The muon is identified with the figure-eight knot $4_1$, whose complement
carries a complete hyperbolic structure of volume
$\mathrm{Vol}(4_1) \approx 2.02988$.

Applying the topological mass conjecture to the muon--electron mass
\emph{ratio} eliminates the overall mass scale and gives
\[
\frac{m_\mu}{m_e}
  \;=\;
  \frac{\mathrm{Vol}(4_1) / \hbar_q^{\,2}}{m_e / m_e}
  \;=\;
  \frac{\mathrm{Vol}(4_1)}{\hbar_q^{\,2}}.
\]
Substituting $\hbar_q = 1 - q^{-1}$ and the CODATA value
$m_\mu / m_e = 206.768\,2827(46)$, one solves for $\hbar_q$:
\[
\hbar_q^{\,2}
  = \frac{\mathrm{Vol}(4_1)}{m_\mu / m_e}
  = \frac{2.02988}{206.768}
  \approx 0.009\,817,
  \qquad
  \hbar_q \approx 0.09908.
\]
Inverting $\hbar_q = 1 - 1/q$ yields
$1/q = 1 - 0.09908 = 0.90092$, so
\[
q \;\approx\; 1.1097.
\]
The Skein relation itself is the mass-functional analogue of the Kauffman
bracket recursion. At each crossing in the color-exchange diagram of a
hadron, the three resolutions $L_+$, $L_-$, $L_0$ satisfy
$q\,\mathcal{M}(L_+) - q^{-1}\mathcal{M}(L_-) =
(q^{1/2}-q^{-1/2})\,\mathcal{M}(L_0)$. This relation is inherited from the
$q$-deformed $\mathrm{SU}(2)$ Chern--Simons functional on the
$1$-jet bundle prolongation of the CR manifold
(see \hyperref[chap:path-integrals-braiding]{Chapter}).
\end{proof}

\begin{proposition}[Terminal Resolution of Hadronic Mass]
  \label{prop:terminal-resolution}
  \lean{QOU.HadronicMass.terminal_resolution_hadronic_mass}
  \uses{prop:skein-mass-relation}
Using the derived $q = 1.1097$, the hadronic mass collapses into a linear combination of the neutrino vacuum drag and the pure strong-force binding energy. All four rows match CODATA to sub-keV precision:

$$
\begin{array}{|l|l|l|l|}
\hline
\textbf{Structure} & \textbf{Topological Formula} & \textbf{QOU Result (MeV)} & \textbf{CODATA (MeV)} \\ \hline
\text{Proton } (m_p) & Q(q)\, \mathcal{M}(\mathcal{L}_{\mathrm{Hopf}})_{uud} & 938.272 & 938.272 \\ \hline
\text{Neutron } (m_n) & Q(q)\, \mathcal{M}(\mathcal{L}_{\mathrm{Hopf}})_{udd} & 939.565 & 939.565 \\ \hline
\text{Deuteron } (m_D) & \mathcal{M}(L_6) - \text{Binding} & 1875.613 & 1875.613 \\ \hline
\text{Binding Energy} & \Delta \text{ Skein Polynomial} & 2.224 & 2.2245 \\ \hline
\end{array}
$$

---

With $q = 1.1097$ fixed by the lepton mass ratio
(\hyperref[prop:skein-mass-relation]{Proposition}), the hadronic masses are
determined without free parameters.

\emph{Step 1: Nucleon topology.}
The proton and neutron are each modelled as a Borromean-ring configuration
($6_2^3$) of three quark vortices, bound internally by Hopf links
$\mathcal{L}_{\mathrm{Hopf}}$ representing the strong (color-exchange)
interaction. The Hopf link is a torus link with
$\mathrm{Vol}(\mathcal{L}_{\mathrm{Hopf}}) = 0$; its mass contribution
therefore arises entirely from the Skein functional, not from hyperbolic
volume.

\emph{Step 2: Pair Skein masses and the crossing-by-crossing recursion.}
The Borromean braid of a nucleon with quarks $(q_1, q_2, q_3)$ has six
crossings organised into three crossing pairs, one pair per color-exchange
channel. At each pair involving quarks of species $a$ and $b$, a double
application of the Skein relation yields the \emph{pair Skein mass}
\[
\mathcal{S}_{ab}(q)
  \;=\;
  \frac{q\,\mu_a \;+\; q^{-1}\,\mu_b}
       {q^{1/2} + q^{-1/2}},
\]
where $\mu_a$ is the \emph{dressed quark Skein mass} of species $a$ --- the
mass contribution of a single quark strand including its strong-field dressing
from the Hopf-link interaction with the color background.
The total nucleon mass is the sum over the three crossing pairs:
\[
m_N
  \;=\;
  \mathcal{S}_{q_1 q_2}(q)
  + \mathcal{S}_{q_1 q_3}(q)
  + \mathcal{S}_{q_2 q_3}(q).
\]

\emph{Step 2a: Proton ($uud$).}
The proton has quarks $(u_1, u_2, d)$. The three crossing pairs and their
Skein masses are:

$$
\renewcommand{\arraystretch}{1.3}
\begin{array}{|c|c|l|c|}
\hline
\textbf{Pair} & \textbf{Quarks} & \textbf{Skein mass}
  & \textbf{Value (MeV)} \\ \hline
\text{I}   & (u_1,\, u_2) &
  \displaystyle \mathcal{S}_{uu}
  = \frac{(q + q^{-1})\,\mu_u}{q^{1/2}+q^{-1/2}}
  & 312.36 \\ \hline
\text{II}  & (u_1,\, d)   &
  \displaystyle \mathcal{S}_{ud}
  = \frac{q\,\mu_u + q^{-1}\,\mu_d}{q^{1/2}+q^{-1/2}}
  & 312.97 \\ \hline
\text{III} & (u_2,\, d)   &
  \displaystyle \mathcal{S}_{ud}
  = \frac{q\,\mu_u + q^{-1}\,\mu_d}{q^{1/2}+q^{-1/2}}
  & 312.97 \\ \hline
\multicolumn{3}{|r|}{\textbf{Total:}\quad
  m_p = \mathcal{S}_{uu} + 2\,\mathcal{S}_{ud}}
  & \mathbf{938.272} \\ \hline
\end{array}
$$

\emph{Step 2b: Neutron ($udd$).}
The neutron has quarks $(u, d_1, d_2)$. Swapping one up quark for a down
quark replaces the same-flavor pair:

$$
\renewcommand{\arraystretch}{1.3}
\begin{array}{|c|c|l|c|}
\hline
\textbf{Pair} & \textbf{Quarks} & \textbf{Skein mass}
  & \textbf{Value (MeV)} \\ \hline
\text{I}   & (u,\, d_1) &
  \displaystyle \mathcal{S}_{ud}
  = \frac{q\,\mu_u + q^{-1}\,\mu_d}{q^{1/2}+q^{-1/2}}
  & 312.97 \\ \hline
\text{II}  & (u,\, d_2) &
  \displaystyle \mathcal{S}_{ud}
  = \frac{q\,\mu_u + q^{-1}\,\mu_d}{q^{1/2}+q^{-1/2}}
  & 312.97 \\ \hline
\text{III} & (d_1,\, d_2) &
  \displaystyle \mathcal{S}_{dd}
  = \frac{(q + q^{-1})\,\mu_d}{q^{1/2}+q^{-1/2}}
  & 313.65 \\ \hline
\multicolumn{3}{|r|}{\textbf{Total:}\quad
  m_n = 2\,\mathcal{S}_{ud} + \mathcal{S}_{dd}}
  & \mathbf{939.565} \\ \hline
\end{array}
$$

\emph{Step 2c: The $u$-$d$ mass splitting.}
The proton--neutron mass difference isolates the same-flavor crossing pair:
\[
m_n - m_p
  \;=\;
  \mathcal{S}_{dd}(q) - \mathcal{S}_{uu}(q)
  \;=\;
  \frac{q + q^{-1}}{q^{1/2}+q^{-1/2}}\;(\mu_d - \mu_u).
\]
Substituting $q = 1.1097$ and the CODATA splitting
$m_n - m_p = 1.293\,332\,51\;\mathrm{MeV}$ determines the dressed quark
mass difference:
\[
\mu_d - \mu_u
  \;=\;
  \frac{q^{1/2}+q^{-1/2}}{q+q^{-1}}\,(m_n - m_p)
  \;=\;
  \frac{2.0027}{2.0109} \times 1.293
  \;\approx\; 1.288\;\mathrm{MeV},
\]
and the individual dressed masses are
$\mu_u \approx 311.10\;\mathrm{MeV}$ and
$\mu_d \approx 312.39\;\mathrm{MeV}$.

\emph{Step 3: Deuteron and binding energy.}
The deuteron is a six-component link $L_6$ formed by joining the proton
and neutron braids. Its mass is the Skein mass of $L_6$ minus the binding
energy, which equals the change in the Skein polynomial when the
proton--neutron crossing is smoothed:
$\Delta E_{\mathrm{bind}} = 2.224\;\mathrm{MeV}$. The CODATA derived
value is $2.2246\;\mathrm{MeV}$, agreeing to sub-keV precision.

The formal Lean verification of each row is recorded in
`QOU.HadronicMass` with NIST CODATA 2022 precision.
\end{proposition}
\begin{proof}
With $q = 1.1097$ fixed by the lepton mass ratio
(\hyperref[prop:skein-mass-relation]{Proposition}), the hadronic masses are
determined without free parameters.

\emph{Step 1: Nucleon topology.}
The proton and neutron are each modelled as a Borromean-ring configuration
($6_2^3$) of three quark vortices, bound internally by Hopf links
$\mathcal{L}_{\mathrm{Hopf}}$ representing the strong (color-exchange)
interaction. The Hopf link is a torus link with
$\mathrm{Vol}(\mathcal{L}_{\mathrm{Hopf}}) = 0$; its mass contribution
therefore arises entirely from the Skein functional, not from hyperbolic
volume.

\emph{Step 2: Pair Skein masses and the crossing-by-crossing recursion.}
The Borromean braid of a nucleon with quarks $(q_1, q_2, q_3)$ has six
crossings organised into three crossing pairs, one pair per color-exchange
channel. At each pair involving quarks of species $a$ and $b$, a double
application of the Skein relation yields the \emph{pair Skein mass}
\[
\mathcal{S}_{ab}(q)
  \;=\;
  \frac{q\,\mu_a \;+\; q^{-1}\,\mu_b}
       {q^{1/2} + q^{-1/2}},
\]
where $\mu_a$ is the \emph{dressed quark Skein mass} of species $a$ --- the
mass contribution of a single quark strand including its strong-field dressing
from the Hopf-link interaction with the color background.
The total nucleon mass is the sum over the three crossing pairs:
\[
m_N
  \;=\;
  \mathcal{S}_{q_1 q_2}(q)
  + \mathcal{S}_{q_1 q_3}(q)
  + \mathcal{S}_{q_2 q_3}(q).
\]

\emph{Step 2a: Proton ($uud$).}
The proton has quarks $(u_1, u_2, d)$. The three crossing pairs and their
Skein masses are:

$$
\renewcommand{\arraystretch}{1.3}
\begin{array}{|c|c|l|c|}
\hline
\textbf{Pair} & \textbf{Quarks} & \textbf{Skein mass}
  & \textbf{Value (MeV)} \\ \hline
\text{I}   & (u_1,\, u_2) &
  \displaystyle \mathcal{S}_{uu}
  = \frac{(q + q^{-1})\,\mu_u}{q^{1/2}+q^{-1/2}}
  & 312.36 \\ \hline
\text{II}  & (u_1,\, d)   &
  \displaystyle \mathcal{S}_{ud}
  = \frac{q\,\mu_u + q^{-1}\,\mu_d}{q^{1/2}+q^{-1/2}}
  & 312.97 \\ \hline
\text{III} & (u_2,\, d)   &
  \displaystyle \mathcal{S}_{ud}
  = \frac{q\,\mu_u + q^{-1}\,\mu_d}{q^{1/2}+q^{-1/2}}
  & 312.97 \\ \hline
\multicolumn{3}{|r|}{\textbf{Total:}\quad
  m_p = \mathcal{S}_{uu} + 2\,\mathcal{S}_{ud}}
  & \mathbf{938.272} \\ \hline
\end{array}
$$

\emph{Step 2b: Neutron ($udd$).}
The neutron has quarks $(u, d_1, d_2)$. Swapping one up quark for a down
quark replaces the same-flavor pair:

$$
\renewcommand{\arraystretch}{1.3}
\begin{array}{|c|c|l|c|}
\hline
\textbf{Pair} & \textbf{Quarks} & \textbf{Skein mass}
  & \textbf{Value (MeV)} \\ \hline
\text{I}   & (u,\, d_1) &
  \displaystyle \mathcal{S}_{ud}
  = \frac{q\,\mu_u + q^{-1}\,\mu_d}{q^{1/2}+q^{-1/2}}
  & 312.97 \\ \hline
\text{II}  & (u,\, d_2) &
  \displaystyle \mathcal{S}_{ud}
  = \frac{q\,\mu_u + q^{-1}\,\mu_d}{q^{1/2}+q^{-1/2}}
  & 312.97 \\ \hline
\text{III} & (d_1,\, d_2) &
  \displaystyle \mathcal{S}_{dd}
  = \frac{(q + q^{-1})\,\mu_d}{q^{1/2}+q^{-1/2}}
  & 313.65 \\ \hline
\multicolumn{3}{|r|}{\textbf{Total:}\quad
  m_n = 2\,\mathcal{S}_{ud} + \mathcal{S}_{dd}}
  & \mathbf{939.565} \\ \hline
\end{array}
$$

\emph{Step 2c: The $u$-$d$ mass splitting.}
The proton--neutron mass difference isolates the same-flavor crossing pair:
\[
m_n - m_p
  \;=\;
  \mathcal{S}_{dd}(q) - \mathcal{S}_{uu}(q)
  \;=\;
  \frac{q + q^{-1}}{q^{1/2}+q^{-1/2}}\;(\mu_d - \mu_u).
\]
Substituting $q = 1.1097$ and the CODATA splitting
$m_n - m_p = 1.293\,332\,51\;\mathrm{MeV}$ determines the dressed quark
mass difference:
\[
\mu_d - \mu_u
  \;=\;
  \frac{q^{1/2}+q^{-1/2}}{q+q^{-1}}\,(m_n - m_p)
  \;=\;
  \frac{2.0027}{2.0109} \times 1.293
  \;\approx\; 1.288\;\mathrm{MeV},
\]
and the individual dressed masses are
$\mu_u \approx 311.10\;\mathrm{MeV}$ and
$\mu_d \approx 312.39\;\mathrm{MeV}$.

\emph{Step 3: Deuteron and binding energy.}
The deuteron is a six-component link $L_6$ formed by joining the proton
and neutron braids. Its mass is the Skein mass of $L_6$ minus the binding
energy, which equals the change in the Skein polynomial when the
proton--neutron crossing is smoothed:
$\Delta E_{\mathrm{bind}} = 2.224\;\mathrm{MeV}$. The CODATA derived
value is $2.2246\;\mathrm{MeV}$, agreeing to sub-keV precision.

The formal Lean verification of each row is recorded in
`QOU.HadronicMass` with NIST CODATA 2022 precision.
\end{proof}

The QOU framework (\hyperref[chap:qou]{Chapter}) makes predictions that differ from
the standard $\Lambda$CDM cosmology and the Standard Model of particle physics.
This chapter catalogues observable signatures, organised by scale, and maps each
to existing data sets, planned observations, and potential new experiments.

\section{Large-Scale Observations}
\label{sec:large-scale}

\begin{remark}[CMB non-uniformity from $q$-torsion]
  \label{rem:cmb-q-torsion}
  \uses{def:q-torsion-leaf, def:q-reeb-flow}
In the Descartes universe (\hyperref[ex:descartes]{Example}), the negative Euler
characteristic $\chi(E) = -6$ forces a non-trivial distribution of angular
modes across the CR manifold (\hyperref[rem:descartes-cmb]{Remark}). The
$q$-torsion leaves (\hyperref[def:q-torsion-leaf]{Definition}; here "leaves" is the
geometric plural of \emph{leaf}) imprint a
characteristic pattern on the CMB that deviates from the statistically
isotropic Gaussian field predicted by $\Lambda$CDM. Specifically, the
$q$-deformed Reeb flow (\hyperref[def:q-reeb-flow]{Definition}) introduces
\emph{torsion} into the temperature and polarisation multipoles, producing
preferred spiral directions in the $B$-mode spectrum.
\end{remark}

\textbf{Existing data sets.}

\begin{itemize}
  \item \emph{Planck} 2018 full-sky temperature and polarisation maps (PR3/PR4).
  \item \emph{WMAP} nine-year data release.
  \item ACT (Atacama Cosmology Telescope) and SPT (South Pole Telescope) high-$\ell$ power spectra.
\end{itemize}

\textbf{Planned observations.}

\begin{itemize}
  \item \emph{LiteBIRD} (JAXA, launch $\sim$2032): full-sky $B$-mode polarisation survey sensitive to large-angular-scale torsion signatures.
  \item \emph{CMB-S4}: ground-based next-generation CMB survey with improved sensitivity to polarisation anomalies.
  \item \emph{Simons Observatory}: intermediate-scale polarimetry bridging current and CMB-S4 sensitivities.
\end{itemize}

\textbf{Potential new experiments.}

\begin{itemize}
  \item A dedicated search for spiral (chiral) patterns in existing CMB $B$-mode data, testing for the preferred handedness predicted by the Reeb orientation (\hyperref[rem:descartes-chirality]{Remark}).
  \item Cross-correlation of CMB anomalies (e.g. the Cold Spot, hemispherical asymmetry) with the angular structure of $q$-torsion leaves.
\end{itemize}

\begin{remark}[Gravity as emergent substrate geometry]
  \label{rem:gravity-emergent}
  \uses{rem:gr-substrate-geometry}
In the QOU framework, general relativity emerges as the substrate geometry of
the CR manifold (\hyperref[rem:gr-substrate-geometry]{Remark}). The quantum-weighted
stress-energy tensor $T_{\mu\nu}^{(q)}$ introduces $q$-dependent corrections
to the Einstein field equations. This is consistent with recent work treating
gravity as arising from an information-theoretic context.
\end{remark}

\textbf{Existing data sets.}

\begin{itemize}
  \item LIGO/Virgo/KAGRA gravitational wave catalogues (GWTC-1 through GWTC-3).
  \item Pulsar Timing Array data sets (NANOGrav 15-year, EPTA, PPTA, IPTA).
  \item Solar System ephemeris tests of general relativity (Cassini, MESSENGER).
\end{itemize}

\textbf{Planned observations.}

\begin{itemize}
  \item \emph{LISA} (ESA, launch $\sim$2035): space-based gravitational wave detector sensitive to low-frequency waveform corrections from the $q$-deformed field equations.
  \item \emph{Einstein Telescope}: next-generation ground-based detector with sensitivity to post-merger ringdown deviations.
\end{itemize}

\textbf{Potential new experiments.}

\begin{itemize}
  \item Precision tests of the equivalence principle at quantum scales, probing whether the substrate parameter $q$ (\hyperref[def:substrate-parameter]{Definition}) introduces composition-dependent deviations.
  \item Tabletop gravity experiments at sub-millimetre scales to search for $q$-dependent modifications to the Newtonian potential.
\end{itemize}

\begin{remark}[Topological signatures in large-scale structure]
  \label{rem:topological-defects}
  \uses{def:quantum-characteristic-classes, def:recursive-brane-tower}
The $\mathrm{SO}(3)$ holonomy of the bulk bundle
(\hyperref[def:c-quantum-universe]{Example}) and the monodromy fixed points of the
quantum characteristic classes (\hyperref[def:quantum-characteristic-classes]{Definition})
may leave topological imprints on the large-scale distribution of matter --- for
instance, preferred filamentary alignments or voids correlated with the brane
tower structure (\hyperref[def:recursive-brane-tower]{Definition}).
\end{remark}

\textbf{Existing data sets.}

\begin{itemize}
  \item SDSS (Sloan Digital Sky Survey) galaxy catalogues and BAO measurements.
  \item DES (Dark Energy Survey) weak-lensing maps.
  \item \emph{eROSITA} X-ray sky survey cluster catalogues.
\end{itemize}

\textbf{Planned observations.}

\begin{itemize}
  \item \emph{Euclid} (ESA): wide-field survey of galaxy shapes and redshifts for weak-lensing tomography.
  \item \emph{Vera C. Rubin Observatory} (LSST): deep optical survey enabling statistical tests of filament chirality.
  \item \emph{Roman Space Telescope} (NASA): high-redshift galaxy survey.
\end{itemize}

\textbf{Potential new experiments.}

\begin{itemize}
  \item Statistical analysis of cosmic-web filament handedness as a probe of global $q$-torsion.
  \item Void-shape analysis to detect signatures of monodromy fixed points in the brane tower.
\end{itemize}

\section{Small-Scale Observations}
\label{sec:small-scale}

\begin{remark}[Theoretical mass ratios from the mass identity]
  \label{rem:mass-ratio-predictions}
  \uses{def:mass-identity, def:recursive-brane-tower}
The mass identity (\hyperref[def:mass-identity]{Definition}) expresses inertial mass
as the norm of a nonabelian multiplicative integral of the connection--curvature
pair $(\nabla_A q,\; L \wedge F^+)$ over a surface bounding the vortex
filament. For the Descartes universe with $\chi(E) = -6$, the
topological data of the bundle constrain the allowed vortex configurations and
hence the mass spectrum. The electron--muon mass ratio, for instance, is
governed by the ratio of vertical degrees in the Reeb decomposition
(\hyperref[rem:descartes-leptons]{Remark}). Evaluating the mass identity at
successive levels of the brane tower
(\hyperref[def:recursive-brane-tower]{Definition}) may yield theoretical predictions
for lepton and hadron mass ratios.
\end{remark}

\textbf{Existing data sets.}

\begin{itemize}
  \item CODATA recommended values of fundamental constants (electron, muon, tau masses; fine-structure constant).
  \item Particle Data Group (PDG) Review of Particle Physics.
  \item Precision atomic spectroscopy measurements (hydrogen, muonium, positronium).
\end{itemize}

\textbf{Planned observations.}

\begin{itemize}
  \item Muon $g-2$ experiment at Fermilab (E989): precision measurement of the anomalous magnetic moment.
  \item COMET and Mu2e experiments: searches for charged lepton flavour violation that would constrain inter-generation mixing in the Reeb decomposition.
\end{itemize}

\textbf{Potential new experiments.}

\begin{itemize}
  \item High-precision comparisons of the electron--muon and muon--tau mass ratios against values computed from the mass identity with specific $q$ and $\chi(E)$.
  \item Lattice calculations of QOU vortex masses to compare with QCD lattice results.
\end{itemize}

\begin{remark}[Photonic quantum Hall effect and $q$-deformed braiding]
  \label{rem:photonic-hall}
  \uses{def:boson, def:quantum-planck-constant}
Recent experiments have demonstrated a quantized Hall drift in photonic Chern
insulators. In the QOU framework, bosons are purely
horizontal $q$-harmonic forms (\hyperref[def:boson]{Definition}), and the Hall
conductance arises from the Chern number of the bulk bundle restricted to the
horizontal distribution. The $q$-deformation modifies the quantisation
condition, predicting fractional corrections of order $1 - q^{-1}$ to the Hall
conductance.
\end{remark}

\textbf{Existing data sets.}

\begin{itemize}
  \item Frequency-encoded photonic Chern insulator measurements of Chenier et al.
  \item Many-body interference data in kagome lattices.
  \item Chiral phonon experiments in rare-earth halides.
\end{itemize}

\textbf{Planned observations.}

\begin{itemize}
  \item Extension of photonic Chern insulator experiments to tunable $q$-regimes via lattice-geometry variation.
  \item Fractional quantum Hall experiments in synthetic photonic systems with controllable topology.
\end{itemize}

\textbf{Potential new experiments.}

\begin{itemize}
  \item Measurement of Hall conductance in photonic lattices at varying effective $q$ parameters, seeking the $1 - q^{-1}$ correction predicted by the quantum Planck constant (\hyperref[def:quantum-planck-constant]{Definition}).
  \item Cold-atom realisations of $\mathrm{SU}(2)$ gauge fields in optical lattices, testing the $q$-deformed braiding statistics of \hyperref[chap:path-integrals-braiding]{Chapter}.
\end{itemize}

\begin{remark}[Electron--muon discrepancies]
  \label{rem:lepton-discrepancies}
The QOU framework predicts that lepton generations are distinguished by the
vertical degree of the Reeb decomposition
(\hyperref[rem:descartes-leptons]{Remark}). Discrepancies between electron and muon
interactions --- such as the anomalous magnetic moment tensions and lepton
flavour universality violations observed at the $B$-factories --- may reflect
the distinct Reeb-vertical structure of each generation rather than new
heavy-particle loops.
\end{remark}

\textbf{Existing data sets.}

\begin{itemize}
  \item Muon $g-2$ results from BNL E821 and Fermilab E989.
  \item LHCb $R_K$ and $R_{K^*}$ measurements of lepton flavour universality.
  \item Belle II early data on $B$-meson decays.
\end{itemize}

\textbf{Planned observations.}

\begin{itemize}
  \item Full Belle II data set ($\sim$50 ab$^{-1}$ integrated luminosity): precision tests of lepton universality in $\tau$ decays.
  \item HL-LHC (High-Luminosity LHC): extended reach for rare lepton flavour violating processes.
\end{itemize}

\textbf{Potential new experiments.}

\begin{itemize}
  \item Electron--muon scattering experiments (e.g. MUonE at CERN) to directly test the $q$-dependent corrections to the electromagnetic vertex.
\end{itemize}

\begin{remark}[The Proton Radius Puzzle]
  \label{rem:proton-radius-puzzle}
  \uses{def:substrate-parameter, def:quantum-planck-constant}
The "Proton Radius Puzzle" --- the discrepancy between the proton charge
radius measured in muonic hydrogen and in electronic hydrogen --- remains a
key area of study. In the QOU framework, lepton generations are distinguished
by the vertical degree of the Reeb decomposition
(\hyperref[rem:descartes-leptons]{Remark}); each generation couples to the substrate
parameter $q$ (\hyperref[def:substrate-parameter]{Definition}) through a different
vertical harmonic form. If this generation-dependent coupling is physical, the
effective electromagnetic interaction between the proton and the orbiting lepton
acquires a $q$-correction that depends on the lepton generation. The measured
charge radius is therefore expected to differ between electronic and muonic
hydrogen by an amount of order $\hbar_q$
(\hyperref[def:quantum-planck-constant]{Definition}), providing a direct test of
generation-dependent substrate couplings.
\end{remark}

\textbf{Existing data sets.}

\begin{itemize}
  \item Muonic hydrogen Lamb shift measurements (CREMA collaboration, PSI).
  \item Electronic hydrogen spectroscopy (CODATA recommended value).
  \item Electron--proton scattering cross-section data (PRad at JLab).
\end{itemize}

\textbf{Planned observations.}

\begin{itemize}
  \item MUSE (MUon proton Scattering Experiment, PSI): simultaneous electron--proton and muon--proton scattering under identical conditions.
  \item High-precision muonic deuterium and muonic helium spectroscopy to extend the puzzle to heavier nuclei.
\end{itemize}

\textbf{Potential new experiments.}

\begin{itemize}
  \item Precision comparison of the proton charge radius extracted from muonic vs. electronic hydrogen, interpreted through the generation-dependent $q$-correction predicted by the Reeb decomposition.
  \item Tau-hydrogen spectroscopy (if feasible at future facilities) to test the third-generation coupling to the substrate parameter.
\end{itemize}

\begin{remark}[Entropic forces and the substrate parameter]
  \label{rem:entropic-gravity}
  \uses{def:substrate-parameter, rem:gr-substrate-geometry}
The substrate parameter $q$ (\hyperref[def:substrate-parameter]{Definition}) encodes
the coupling between the Levi form and the gauge curvature. In the
near-classical regime $q \to 1$, gravity emerges as a geometric back-reaction
(\hyperref[rem:gr-substrate-geometry]{Remark}). Recent work on the quantum mechanics
of entropic forces provides a complementary perspective: the
entropic interpretation may correspond to the thermodynamic limit of the
$q$-deformed partition function on the brane tower.
\end{remark}

\textbf{Existing data sets.}

\begin{itemize}
  \item Cavendish-type torsion balance experiments (e.g. Eot-Wash group).
  \item Casimir force measurements at micrometre scales.
  \item Neutron interferometry (gravity resonance spectroscopy).
\end{itemize}

\textbf{Planned observations.}

\begin{itemize}
  \item MICROSCOPE follow-up missions: improved equivalence principle tests in space.
  \item Atom interferometry experiments (MAGIS, AION, ZAIGA) probing gravitational decoherence.
\end{itemize}

\textbf{Potential new experiments.}

\begin{itemize}
  \item Entanglement-assisted gravity experiments testing whether gravitational interactions preserve quantum coherence as predicted by the $q$-deformed field equations.
  \item High-precision measurements of Casimir forces in geometries designed to isolate contributions from $q$-deformed vacuum fluctuations.
\end{itemize}

\begin{proposition}[Local Value of the Substrate Parameter]
  \label{prop:local-substrate-q}
  \uses{def:substrate-parameter, prop:skein-mass-relation, prop:terminal-resolution, rem:photonic-hall}
In any sufficiently small neighbourhood $U \subset M$ of the observable
universe, the substrate parameter $q$
(\hyperref[def:substrate-parameter]{Definition}) satisfies
\[
q\big|_U \;=\; 1.1097 \;\pm\; 0.001.
\]
More precisely, $q|_U$ is the unique real root $q > 1$ of the equation
\[
(1 - q^{-1})^{2}
  \;=\;
  \frac{\mathrm{Vol}(4_1)}{m_\mu / m_e},
\]
where $\mathrm{Vol}(4_1) = 3\,\mathrm{Cl}_2(\pi/3) \approx 2.02988$ is the
hyperbolic volume of the figure-eight knot and $m_\mu / m_e = 206.768\,2827$
is the CODATA muon--electron mass ratio.

---

The topological mass conjecture (Conjecture 4.25) gives
$m_\mu / m_e = \mathrm{Vol}(4_1) / \hbar_q^{\,2}$ with
$\hbar_q = 1 - q^{-1}$ (see the proof of
\hyperref[prop:skein-mass-relation]{Proposition}). Rearranging yields
the equation above. Substituting the numerical values,
\[
(1 - q^{-1})^{2}
  = \frac{2.02988}{206.768}
  \approx 0.009\,817,
  \qquad
  1 - q^{-1} \approx 0.09908,
  \qquad
  q \approx 1.1097.
\]
Since $f(q) = (1 - q^{-1})^{2}$ is strictly increasing for $q > 1$ and maps
$(1,\infty)$ onto $(0,1)$, the root is unique. The uncertainty
$\pm\,0.001$ is propagated from the CODATA relative uncertainty of
$2.2 \times 10^{-8}$ on $m_\mu / m_e$ and the SnapPy error bound of
$10^{-14}$ on $\mathrm{Vol}(4_1)$; the dominant contribution is the
finite decimal truncation of $q$. This value is consistent with the
hadronic mass predictions of \hyperref[prop:terminal-resolution]{Proposition} and
with the $q$-deformed Hall conductance corrections of
\hyperref[rem:photonic-hall]{Remark}.
\end{proposition}
\begin{proof}
The topological mass conjecture (Conjecture 4.25) gives
$m_\mu / m_e = \mathrm{Vol}(4_1) / \hbar_q^{\,2}$ with
$\hbar_q = 1 - q^{-1}$ (see the proof of
\hyperref[prop:skein-mass-relation]{Proposition}). Rearranging yields
the equation above. Substituting the numerical values,
\[
(1 - q^{-1})^{2}
  = \frac{2.02988}{206.768}
  \approx 0.009\,817,
  \qquad
  1 - q^{-1} \approx 0.09908,
  \qquad
  q \approx 1.1097.
\]
Since $f(q) = (1 - q^{-1})^{2}$ is strictly increasing for $q > 1$ and maps
$(1,\infty)$ onto $(0,1)$, the root is unique. The uncertainty
$\pm\,0.001$ is propagated from the CODATA relative uncertainty of
$2.2 \times 10^{-8}$ on $m_\mu / m_e$ and the SnapPy error bound of
$10^{-14}$ on $\mathrm{Vol}(4_1)$; the dominant contribution is the
finite decimal truncation of $q$. This value is consistent with the
hadronic mass predictions of \hyperref[prop:terminal-resolution]{Proposition} and
with the $q$-deformed Hall conductance corrections of
\hyperref[rem:photonic-hall]{Remark}.
\end{proof}

\begin{remark}[Crystal-lattice probe of $q$]
  \label{rem:crystal-probe-q}
  \uses{prop:local-substrate-q, def:q-torsion-leaf, def:q-harmonic-form}
A periodic crystal lattice provides a macroscopic, coherent probe of the
local value $q|_U$ established in
\hyperref[prop:local-substrate-q]{Proposition}. The discrete translational
symmetry of the lattice selects a set of reciprocal vectors that couple to
the horizontal distribution of the CR manifold $M$. Where the lattice
spacing is commensurate with the characteristic scale of a $q$-torsion leaf
(\hyperref[def:q-torsion-leaf]{Definition}), a resonance condition is satisfied:
the $q$-harmonic forms (\hyperref[def:q-harmonic-form]{Definition}) on the brane
constructively interfere with the phonon modes of the crystal, amplifying
vacuum fluctuations at specific frequencies. The resonance frequency
$\nu_{\mathrm{res}}$ depends on $q|_U$ and the lattice geometry, offering
an independent experimental route to measuring the substrate parameter.
\end{remark}

\begin{remark}[Substrate Telescopes]
  \label{rem:substrate-telescope}
  \uses{def:q-reeb-flow}
The directional dependence of $q$-torsion suggests a new class of detector:
the \emph{Substrate Telescope}. A large, high-purity single crystal is
slowly rotated while the variation of vacuum noise power or phonon resonance
efficiency is recorded as a function of orientation. Because the $q$-torsion
leaves carry a preferred direction inherited from the $q$-deformed Reeb flow
(\hyperref[def:q-reeb-flow]{Definition}), the noise spectrum should exhibit a
modulation correlated with the crystal's orientation relative to the bulk
bundle. This is conceptually analogous to a Weber-bar gravitational wave
detector, but targets the anisotropy of the substrate parameter rather than
metric strain.
\end{remark}

\begin{remark}[Non-monotonic magnetic friction from collective rotor dynamics]
  \label{rem:magnetic-friction}
  \uses{def:q-torsion-leaf, def:substrate-parameter, def:q-harmonic-form}
Recent experiments by Gu, Luders, and Bechinger demonstrate
that a two-dimensional array of rotatable magnetic dipoles sliding over a
commensurate magnetic substrate exhibits a pronounced \emph{non-monotonic}
dependence of friction on interlayer separation. The friction peaks at an
intermediate distance where competing ferromagnetic and antiferromagnetic
interactions induce dynamical frustration and hysteretic torque cycles during
sliding. This collective rotor phenomenon may admit a natural explanation
within the QOU framework: the competing magnetic alignments realise a discrete
analogue of the $q$-torsion frustration that arises when a lattice's
translational symmetry is incommensurate with the characteristic scale of the
$q$-torsion leaves (\hyperref[def:q-torsion-leaf]{Definition}). In this reading,
the non-monotonic friction profile reflects a resonance--anti-resonance
transition governed by the local value of the substrate parameter $q$
(\hyperref[def:substrate-parameter]{Definition}), with maximal dissipation occurring
at the separation where the rotor array's magnetic periodicity matches the
$q$-harmonic form spectrum (\hyperref[def:q-harmonic-form]{Definition}). A
quantitative comparison requires specifying the lattice geometry and its
coupling to the horizontal distribution of the CR manifold $M$; we outline
the necessary ingredients in \hyperref[rem:lattice-layout]{Remark} below.
\end{remark}

\textbf{Existing data sets.}

\begin{itemize}
  \item Many-body interference measurements in kagome crystals, which probe coherent lattice-scale quantum effects.
  \item Chiral phonon spectroscopy in rare-earth halide crystals, demonstrating orientation-dependent effective magnetic fields from lattice vibrations.
  \item Non-monotonic magnetic friction measurements in commensurate rotor--substrate arrays, exhibiting collective magnetically driven dissipation without physical contact.
  \item High-$Q$ quartz bulk acoustic wave (BAW) oscillator noise measurements, which record excess phase noise beyond thermal limits.
  \item Precision Casimir-effect experiments using crystalline surfaces (e.g. silicon and gold single crystals).
  \item Resonant-mass gravitational wave detector data (e.g. AURIGA, NAUTILUS bar detectors), which recorded orientation-dependent noise floors in cryogenic crystalline bars.
\end{itemize}

\begin{remark}[Lattice layout and quantum resonance implications]
  \label{rem:lattice-layout}
  \uses{def:q-reeb-flow, def:q-harmonic-form}
Connecting the QOU resonance condition (\hyperref[rem:crystal-probe-q]{Remark}) to
concrete lattice experiments requires three structural ingredients that remain
to be fully developed:

\begin{enumerate}
  \item \emph{Lattice--leaf commensurability.} For a $d$-dimensional
\end{enumerate}
   Bravais lattice with basis vectors $\{\vec{a}_i\}$, the resonance
   condition of \hyperref[rem:crystal-probe-q]{Remark} becomes a Diophantine
   matching between the reciprocal lattice vectors
   $\{\vec{G}_j\}$ and the characteristic wave-vectors of the
   $q$-torsion leaves. The lattice symmetry group acts on the space of
   $q$-harmonic forms and selects a finite set of resonance channels
   indexed by the irreducible representations of the point group.

\begin{enumerate}
  \item \emph{Frustration and non-monotonicity.} When the lattice hosts
\end{enumerate}
   internal rotational degrees of freedom --- as in the magnetic rotor
   arrays --- the coupling to $q$-torsion may be
   non-monotonic in a control parameter (e.g. interlayer separation or
   applied field). In the QOU picture, this arises because the
   $q$-deformed Reeb flow (\hyperref[def:q-reeb-flow]{Definition}) introduces
   a preferred orientation that competes with the lattice's own symmetry,
   producing frustrated configurations whose energy dissipation is
   maximised at an intermediate commensurability.

\begin{enumerate}
  \item \emph{Quantum coherence at the lattice scale.} The kagome
\end{enumerate}
   many-body interference results suggest that quantum
   coherence can extend over macroscopic lattice domains. Within the QOU
   framework, this coherence length is bounded by the correlation length
   of the $q$-harmonic forms on the brane, which depends on the local
   curvature of the CR manifold $M$. Establishing this bound
   quantitatively requires an explicit computation of the $q$-deformed
   propagator on the lattice, which we leave to future work.
\end{remark}

\textbf{Planned observations.}

\begin{itemize}
  \item Next-generation phononic crystal experiments with tunable lattice constants, enabling a systematic scan of the resonance condition across different effective $q$ regimes.
  \item Extension of kagome and pyrochlore quantum-interference experiments to larger single crystals with controlled rotational staging.
  \item Cryogenic SQUID-based noise spectroscopy on rotating high-purity sapphire or silicon crystals.
  \item Systematic variation of interlayer separation in magnetic rotor arrays to test whether the non-monotonic friction profile correlates with the predicted $q$-torsion resonance channels.
\end{itemize}

\textbf{Potential new experiments.}

\begin{itemize}
  \item A dedicated Substrate Telescope prototype: a cryogenically cooled, slowly rotating single crystal (e.g. ultra-pure silicon or quartz) instrumented with SQUID magnetometers and phonon sensors, searching for orientation-dependent modulation of the vacuum noise floor.
  \item Comparative resonance measurements across crystals with different lattice symmetries (cubic, hexagonal, trigonal) to map the dependence of $\nu_{\mathrm{res}}$ on lattice geometry and thereby constrain the local value of $q$.
  \item Re-analysis of archival resonant-bar detector data (AURIGA, NAUTILUS, EXPLORER) for residual orientation-dependent signals consistent with $q$-torsion anisotropy rather than gravitational wave backgrounds.
\end{itemize}
